\documentclass[12pt]{article}
\usepackage[utf8]{inputenc}
\usepackage[T1]{fontenc}
\usepackage[margin=1in]{geometry}
\usepackage{setspace}
\usepackage{xcolor}
\usepackage[colorlinks=true, urlcolor=blue!60!black, linkcolor=blue!60!black]{hyperref}
\setlength{\parindent}{0em}
\setlength{\parskip}{0.6em}
\renewcommand{\baselinestretch}{1.0}
\begin{document}
\begin{center}
{\LARGE \textit{Washington v.\ Davis}, 426 U.S.\ 229 (1976)\par}
\vspace{0.5em}
{\large Case Brief --- Two-Page Edition\par}
\end{center}

\section{Memory Jogger}
Disparate racial impact alone does not violate equal protection; challengers must prove discriminatory purpose (\href{https://supreme.justia.com/cases/federal/us/426/229/}{Justia}).

\section{Detailed Case Facts}
Black applicants to the D.C.\ police challenged ``Test 21,'' a civil-service verbal examination that Black applicants failed at four times the white rate. They alleged no intent to discriminate. The Department had affirmatively recruited Black candidates, 44\% of new recruits were Black, and test scores correlated with training performance.

\section{Procedural History}
The district court granted summary judgment for the officials; the D.C.\ Circuit reversed, applying Title VII's \textit{Griggs} disparate-impact standard to the Constitution. The Supreme Court reversed.

\section{Judicial Votes}
Reversed 7--2. White wrote for Burger, Blackmun, Powell, Rehnquist, and Stevens, with Stewart joining in part; Stevens concurred; Brennan (with Marshall) dissented on the statutory issue.

\section{Holding}
A facially neutral law is not unconstitutional solely because of its racially disproportionate impact; discriminatory purpose is required, though impact is evidence of it. A test may be validated by its link to job-relevant training.

\section{Analysis of Opinions}
\textbf{Majority:} precedent always required purpose, which may be inferred from ``the totality of the relevant facts''; upgrading verbal skills is legitimate; and an impact rule would threaten countless neutral laws. \textbf{Stevens:} objective evidence often proves intent best. \textbf{Brennan:} the test was not shown to measure ``job-specific ability.''

\section{Examples: Future Applications}
\textbf{Same-Side:} a veterans' preference that overwhelmingly favors men (\textit{Feeney}). \textbf{Opposite-Side:} an abrupt rezoning marked by procedural departures and coded remarks (\textit{Arlington Heights}). \textbf{Fence-Sitter:} a bail algorithm using race-correlated proxies adopted despite warnings.

\section{Critique}
Critics such as Charles Lawrence say the intent rule ignores unconscious bias and protects neutral-seeming discrimination; defenders say equal protection targets purposeful denial of equal status and leaves impact remedies to legislatures.

\section{Key Quotations}
\begin{quote}
``Disproportionate impact is not irrelevant, but it is not the sole touchstone of an invidious racial discrimination forbidden by the Constitution. Standing alone, it does not trigger the rule, \textit{McLaughlin v.\ Florida}, 379 U.\ S.\ 184 (1964), that racial classifications are to be subjected to the strictest scrutiny, and are justifiable only by the weightiest of considerations.''
\end{quote}
\noindent --- Justice White, for the Court, 426 U.S.\ at 242.

\end{document}
